\section{Series scope and interface architecture}
\label{sec:series-scope}

\subsection{Structure of the series}

The emergence program is organized as a series of five papers sharing a common substrate
model but developed around logically independent bridges:

\begin{enumerate}
  \item \textbf{Core (this paper).}
  Ontological foundation, continuum model, and linearized wave structure
  (\Cref{sec:continuum-model,sec:waves}).

  \item \textbf{Lorentz and gravity bridge.}
  Imports: branch speeds \eqref{eq:wave-speeds-derived}, transverse--lateral coupling
  \eqref{eq:strain-near-identity}, dual-observer setup.
  Derives: the observer-universality argument for effective Lorentz symmetry at speed $c_T$;
  residual birefringence as a falsifiable signature; quantitative Newtonian reduction of the
  contraction channel sourced by $\partial_i u\,\partial_j u$.

  \item \textbf{Gauge and color bridge.}
  Imports: linearized branch speeds, polarization eigenbundles from the multi-branch wave system,
  the U(3) triplet structure of the three spacelike components.
  Derives: Berry connection $a_\mu = \ii\langle u|\partial_\mu u\rangle$ as gauge potential $A_\mu$;
  Berry curvature as field tensor $F_{\mu\nu}$; Wilczek--Zee connection for the degenerate triplet
  as $\mathfrak{u}(3)$-valued potential; traceless $\mathfrak{su}(3)$ sector and kinematic confinement
  of color-neutral combinations.

  \item \textbf{Matter and mass bridge.}
  Imports: equations of motion \eqref{eq:eom}, geometric nonlinearity, gauge connection from
  Paper~III.
  Derives: soliton eigenproblem as a self-adjoint BVP with nonlinear self-guidance; topological
  identification of stable modes; emergent rest mass as the energy of a self-confined
  closed-loop wave pattern, without invoking the Higgs mechanism.

  \item \textbf{The Bell boundary condition (Paper~V).}
  Imports: substrate determinism, locality, and time-symmetry of the brane action
  (\Cref{subsec:ontology-kinematics}).
  Develops: the uniqueness argument that retrocausality is the only completion consistent with
  this substrate's ontology; causality from soliton chirality; matter and antimatter as
  forward- and backward-propagating worldtubes; entanglement as a V-branching 4D world-volume;
  implications for the simulation program.
\end{enumerate}

\subsection{Interfaces exported by this paper}

The following results, derived in \Cref{sec:continuum-model,sec:waves}, are treated as fixed
interfaces by companion papers.
Companion papers may cite these labels directly without re-derivation.

\begin{itemize}
  \item \textbf{Discrete lattice action:} \Cref{eq:discrete-4d-action} (from
  \Cref{subsec:discrete-4d-action}).
  \item \textbf{Induced metric:} \Cref{eq:induced-metric} and its near-identity expansion
  \Cref{eq:metric-near-identity}.
  \item \textbf{Substrate ontology and kinematics:} \Cref{subsec:ontology-kinematics}
  (imported by Paper~V).
  \item \textbf{Continuum action and constitutive law:}
  \Cref{eq:action,subsec:constitutive-law}, including the geometric quartic coefficient
  \Cref{eq:W4} (imported by Papers~III and~IV).
  \item \textbf{Equations of motion:} \Cref{eq:eom}.
  \item \textbf{Linearized branch speeds:}
  $c_L^2 = k_s a^2/m$ and $c_T^2 = (1-\alpha)\,k_s a^2/m$
  from \Cref{eq:wave-speeds-derived}.
  \item \textbf{Long-wavelength dispersion:} \Cref{eq:linear-dispersion}.
  \item \textbf{Linear wave regime and axis-aligned triplet:} \Cref{sec:waves}
  (imported by Paper~III).
  \item \textbf{In-brane equilibrium and contraction channel:} \Cref{subsec:gravity-channel}
  (imported by Paper~II).
  \item \textbf{Regime of validity:} \Cref{subsec:wave-regime}.
\end{itemize}

\subsection{Scope limits}

To maintain honesty across the series, each companion paper states the non-claims local to its
own bridge --- the limits of what that bridge establishes --- in its introduction.
There is no globally claimed result beyond the per-bridge derivations and the interfaces exported
above.